%%% chapters/06_conclusion.tex — 결론
\section{결론}\label{sec:conclusion}

\subsection{주요 발견}\label{ssec:findings}

본 연구는 LLM 기반 왕복 번역을 활용하여 프로그래밍 언어 간 거리를 정량적으로
측정하는 실험적 방법론을 제안하고, Ollama(qwen2.5-coder:7b)와 OpenAI(gpt-5.4)
두 모델에 대한 비교 실험을 수행하였다.
주요 발견은 다음과 같다.

\begin{enumerate}[leftmargin=2em, label=\arabic*)]
  \item \textbf{모델 특성과 언어 쌍의 상호 작용}:
    두 모델은 언어 쌍에 따라 상반된 강점을 보였다.
    OpenAI는 \textsf{cpp\rarr{}c}에서 100\% 성공하였으나
    \textsf{cpp\rarr{}java/python}에서는 전원 실패하였고,
    Ollama는 세 언어 쌍 모두에서 균형 잡힌 성공률(40--60\%)을 유지하였다.
    이는 모델 규모와 번역 품질의 단순 비례 관계가 성립하지 않음을 보여준다.

  \item \textbf{구조 보존과 의미 보존의 트레이드오프}:
    OpenAI는 의미적 정확성을 달성하면서도 코드 구조를 크게 재구성하는 경향이
    있어(낮은 잔차 유사도, 높은 AST 거리), ``성공''의 질적 측면에서 Ollama와
    차이를 보였다.
    Ollama는 상대적으로 원본에 가까운 구조를 유지하면서 번역하는 보수적 전략을
    취했다.

  \item \textbf{언어 간 거리의 다차원성}:
    잔차 유사도, AST 거리, MOSS 유사도, 수렴 속도는 언어 간 거리의 서로 다른
    측면을 포착하였으며, 이들은 반드시 일치하지 않았다.
    예를 들어, \textsf{cpp\rarr{}python}은 높은 잔차 유사도를 보이면서도
    MOSS 유사도가 $0$\%인 경우가 빈번하여, 토큰 수준의 보존과 구조 수준의 보존이
    독립적일 수 있음을 확인하였다.

  \item \textbf{진동의 발견}:
    Ollama에서 관찰된 2건의 진동(oscillation)은 특정 언어 쌍에서 동등한
    두 가지 코드 표현이 안정적으로 교대할 수 있음을 보여주며, 이는 언어 간
    번역에서 ``정답''이 유일하지 않다는 점을 정량적으로 입증하였다.
\end{enumerate}

\subsection{방법론적 기여}\label{ssec:contributions}

본 연구의 방법론적 기여는 다음과 같다.

\begin{itemize}[leftmargin=2em]
  \item \textbf{자동화된 RTT 파이프라인}:
    코드 번역, 컴파일, 실행 검증, 수렴 판정, 다차원 지표 측정을 완전히
    자동화한 실험 프레임워크를 구현하였다.

  \item \textbf{다차원 거리 지표}:
    단일 지표에 의존하지 않고, 토큰(잔차 유사도)\,$\cdot$\,구조(AST
    거리)\,$\cdot$\,표절 검사(MOSS)\,$\cdot$\,의미(semantic
    pass/fail)\,$\cdot$\,복잡도(lizard 메트릭) 다섯 축을 동시에 측정하여
    언어 간 거리의 복합적 특성을 포착하였다.

  \item \textbf{이종 모델 비교 프로토콜}:
    동일 코퍼스\,$\cdot$\,동일 조건 하에서 이종 LLM의 번역 행동을 체계적으로
    비교할 수 있는 실험 프로토콜을 확립하였다.
\end{itemize}

\subsection{한계 및 향후 연구}\label{ssec:limitations}

본 연구의 주요 한계와 이에 따른 향후 연구 방향은 다음과 같다.

\begin{itemize}[leftmargin=2em]
  \item \textbf{표본 크기}:
    5개 문제\,$\times$\,3개 언어 쌍은 통계적 검정력이 제한적이다.
    표본을 최소 20--30개 문제로 확장해야 언어 쌍별 차이의 통계적 유의성을
    검증할 수 있다.

  \item \textbf{단방향 측정}:
    \cpp{}를 유일한 기준 언어로 사용하여 언어 거리의 비대칭성을 확인할 수
    없었다. 양방향 RTT 실험을 통해
    $D(A \rarr B) \neq D(B \rarr A)$인지 검증할 필요가 있다.

  \item \textbf{모델 다양성}:
    두 모델만 비교하였으므로, Claude, Gemini, CodeLlama 등 다양한 모델로의
    확장을 통해 ``모델 불변적(model-invariant)'' 언어 거리 측정의 가능성을
    탐색해야 한다.

  \item \textbf{재현성 한계}:
    $\text{temperature}=0$에서도 LLM 출력의 완전한 결정론성이 보장되지 않으며,
    모델 버전 업데이트에 따라 결과가 변할 수 있다.
\end{itemize}

이러한 한계에도 불구하고, 본 연구는 LLM 기반 왕복 번역이 프로그래밍 언어 간
거리를 정량적으로 측정하는 유망한 방법론임을 보여주었으며, 이 접근법이
코드 마이그레이션 계획 수립, LLM 번역 능력 벤치마킹, 프로그래밍 언어
계통학(phylogenetics) 연구 등에 활용될 수 있음을 제안한다.
